\providecommand{\SysName}{CARE}

\begin{abstract}
C/C++ refactoring is security-sensitive: edits that simplify code can
silently change cleanup paths, resource ownership, error propagation,
validation checks, or API-visible side effects. Direct LLM-based refactoring is
therefore brittle, because a plausible patch may not apply, may omit relevant
control-flow context, or may weaken resource-lifecycle behavior. This paper
presents \SysName{}, a context-aware automated refactoring engine for
security-preserving C/C++ maintenance. \SysName{} constructs a lightweight
security-aware contextual refactoring graph, detects opportunities such as
resource imbalance, duplicated cleanup, repeated checks, unreachable code, and
security smells, asks an LLM to propose compact patches, and accepts candidates
only after fail-closed validation over patch application, build/test hooks,
resource consistency, semantic stability, and security-regression checks. We
evaluate \SysName{} on OSS50, a benchmark of 50 open-source C/C++ projects. A
full scan covers 49 projects, 25,162 source files, and 266,602 functions,
identifying 142,104 refactoring opportunities. On 487 critical
resource-lifecycle findings, \SysName{} improves patch applicability by
3.1$\times$ and full validation yield by 6.2$\times$ over a direct LLM-only
baseline under the same LLM-call budget, producing 31 fully validated patches
versus 5 for LLM-only. Validator ablations show that resource consistency is
the dominant safety gate: removing it would admit 158 additional patches, 121
of which are classified as false accepts by review. These results suggest that
LLM-based refactoring for systems code benefits from explicit program context
and security-aware validation rather than direct patch prompting alone. The
anonymized implementation and artifacts are available at
\url{https://anonymous.4open.science/r/CARE-03F1}.
\end{abstract}
